\section{Introduction}
\label{sec:introduction}

Enterprise data platforms continuously ingest databases from business units, ERP systems, and exported files. Before analysts, BI tools, or LLM agents can issue meaningful queries, the platform must materialize a \emph{join graph}: which tables join, on which key columns, and in which direction. Text-to-SQL systems over large schemas~\cite{Spider2} and automatic BI modeling~\cite{AutoBI} both depend on this graph. Yet enterprise databases rarely declare foreign keys. Tables arrive as staging copies, snapshots, and exports, and a single business may contain anywhere from 50 to 8\,000 tables. Consequently, the join graph must be discovered from data and metadata, and the discovery process is never finished: new source systems and exports continually arrive, and each new table must be linked before it can participate in queries. We focus on joins within a single business; cross-business joins lie outside the scope of this work.

\introparagraph{Example}
Consider a project-based services firm with \todo{2,010} tables. A project code uniquely identifies each project in a master table, and \todo{k} other tables reference it--ledgers, timesheets, invoices, budgets, together with their staging copies and snapshots. Column names vary (\texttt{PROJECT\_CODE}, \texttt{PRJ\_ID}, \texttt{ENGAGEMENT}); staging copies retain only recent projects; one export pads codes with leading zeros. Simultaneously, dozens of tables contain auto-increment surrogates that begin at 1, status columns whose values lie in $\{1,\dots,5\}$, and a generic column named \texttt{DATA}. A correct join graph links the \todo{k} referencing columns exclusively to the project master and nowhere else. A pairwise method, by contrast, may emit a dense clique among the referencing tables; downstream optimizers or agents then select among equivalent paths arbitrarily and incur the cost of every redundant edge both at discovery time and whenever the graph is placed in a prompt.

\introparagraph{Non-LLM Join Discovery}
Classical approaches decide joins from data statistics. Inclusion-dependency algorithms locate columns whose values are contained in a unique column~\cite{Tschirschnitz17}. Learned detectors score candidates with features such as uniqueness, containment, and name similarity~\cite{Rostin09,ZhangHOPS10,HoPF}; Auto-BI~\cite{AutoBI} further combines a local classifier with global graph optimization. Joinable-table search treats each column as a domain--the set of its distinct values--and retrieves columns whose domains overlap most strongly with a query column. LSH Ensemble~\cite{DBLP:journals/pvldb/ZhuNPM16} established that set containment, rather than Jaccard similarity, is the appropriate overlap measure under skewed domain sizes and indexes it with MinHash sketches at the scale of hundreds of millions of domains. Subsequent systems (JOSIE~\cite{DBLP:conf/sigmod/ZhuDNM19}, MATE~\cite{MATE}, DeepJoin~\cite{2023_dong_deepjoin}, Freyja~\cite{Freyja}) extend the approach to exact overlap, multi-column keys, embeddings, and learned join quality.

These methods are computationally inexpensive and scale to large data lakes, yet they fail on enterprise data in four characteristic ways:
\begin{itemize}[leftmargin=*]
\item \textbf{Accidental overlap produces false joins.} Every surrogate column whose values lie in $1..m$ is contained in every other surrogate whose values lie in $1..n$ ($n\geq m$). Status codes, years, and flags exhibit the same phenomenon. Value statistics alone cannot distinguish these overlaps from genuine references, flooding the graph with spurious edges. LakeBench~\cite{LakeBench} reports analogous failures for joinable-table search. In our databases, \todo{P}\% of inclusion dependencies with containment above 0.9 are not true joins.
\item \textbf{Real joins appear imperfect.} Staging copies, deletions, and snapshots yield only partial containment; formatting differences (leading zeros, prefixes) can drive containment to zero. No single threshold works across an entire database.
\item \textbf{Names carry little reliable semantics.} The same key appears under disparate names, while unrelated keys share generic names such as \texttt{ID}, \texttt{CODE}, or \texttt{DATA}. String similarity cannot resolve either case; the column descriptions and business context that would disambiguate them remain unused.
\item \textbf{Learned models do not transfer.} Detectors and Auto-BI require labeled joins drawn from schemas similar to the target. Each business possesses its own naming conventions and source systems, and obtaining labels is expensive: LakeBench alone required more than 7\,500 human hours. Moreover, joinable-table search returns a ranked list per query column rather than a decisive edge with direction.
\end{itemize}

\introparagraph{LLM-based Join Discovery}
Large language models address these shortcomings because they jointly interpret column names, descriptions, and sample values, and they encode knowledge of common business entities~\cite{dataWrangleNarayan2022}. An LLM can recognize that \texttt{ENGAGEMENT} holds project codes while a surrogate \texttt{ID} does not reference a status table. Current practice therefore submits every pair of tables to an LLM, which decides whether the pair joins and returns the corresponding keys. The approach is accurate, yet its cost scales quadratically with the number of tables: a business of $n$ tables requires $n(n-1)/2$ pair evaluations--approximately 500\,000 for 1\,000 tables and 32 million for 8\,000 tables. Each evaluation itself comprises multiple LLM calls that carry schemas, descriptions, and samples, rendering a single business prohibitively expensive in both tokens and wall-clock time. The cost does not terminate after the initial construction: every newly arriving table must be evaluated against all existing tables, and earlier decisions are never reused. The vast majority of these calls are wasted; only a small fraction of table pairs actually join. Independent pairwise answers may also contradict one another and inject a clique of redundant edges into the graph. Redundant edges are not harmless: downstream tools that select a single path per table pair choose among equivalent paths arbitrarily, so a newly added shortcut can silently alter the tables and keys used by a query; every extra edge is paid for both at discovery time and whenever the graph is embedded in a prompt.

\introparagraph{Hybrid Join Discovery}
A natural response is to combine inexpensive signals with the LLM. Existing hybrids follow two strategies. The first narrows the search space before invoking the LLM. Schuyler~\cite{Schuyler} clusters tables into topics via self-supervised embeddings and observes that joins may then be sought only within and across clusters; our own production pipeline follows a similar idea by tagging tables and grouping them into concepts. Magneto~\cite{Magneto} employs a small model to retrieve candidate column matches for subsequent LLM verification. Topics, however, are not key domains: a customer or project key routinely crosses topic boundaries, so restricting search to a cluster loses true joins while searching across clusters reintroduces a quadratic number of pairs.

The second strategy routes each pair to a cheap proxy or to the LLM. Semantic join operators in LOTUS~\cite{LOTUS} and Palimpzest~\cite{Palimpzest}, together with model cascades~\cite{FrugalGPT}, accept or reject pairs according to embedding similarity and defer only inconclusive cases. SUPG~\cite{SUPG} and BARGAIN~\cite{BARGAIN} set deferral thresholds to guarantee target precision or recall relative to the LLM. Such cascades save little when embeddings fail to encode the join condition; on real data BARGAIN still defers roughly 80\% of pairs~\cite{FDJ}. Featurized-Decomposition Join (FDJ)~\cite{FDJ} reduces cost by up to $10\times$ by extracting features from each record with a linear LLM pass and then deciding joins via inexpensive comparisons on those features.

The third strategy exploits the global structure of the join graph. Auto-BI~\cite{AutoBI} selects edges through a minimum-cost arborescence and therefore assumes a snowflake schema. Nexus~\cite{Nexus} observes that real join graphs are sparse and low-rank because schemas revolve around a few join keys; it recovers the adjacency matrix from metadata alone by low-rank matrix completion and refines candidate scores with entity types inferred by an LLM.

These hybrids lower the per-pair cost yet inherit three fundamental limitations. First, any cheap filter inherits the failure modes of purely statistical methods: a surrogate overlap passes the filter and incurs an LLM call, while a genuine join with zero containment is discarded before the LLM ever sees it. Second, they evaluate a given join predicate over record pairs of two given tables. Join discovery poses a different question--which of thousands of tables join, and on which of their columns. Applied to table or column pairs, these operators judge every pair independently and never exploit the structural regularity that one side of a key is unique, the other side references it, and many tables share a single key domain. Consequently the number of LLM calls continues to grow with the number of candidate pairs. Third, the methods that do exploit global structure use key domains only implicitly. The latent factors of Nexus correspond to key domains, but its output is an undirected matrix of scores without owners or key directions; no step verifies a domain as a unit, the result carries neither a recall guarantee nor a bound on LLM cost, the optimization grows with the number of candidate column pairs, and every arriving table requires completing the whole matrix again.

\begin{table}[t]
\centering
\scriptsize
\setlength{\tabcolsep}{2.5pt}
\setlength{\aboverulesep}{0pt}
\setlength{\belowrulesep}{0pt}
\renewcommand{\arraystretch}{1.35}
\caption{Join key discovery over thousands of tables: non-LLM methods
are cheap but decide from overlap and names, all-pairs LLM prompting
calls the LLM per table pair, hybrids still judge pairs independently;
\ourmethod calls the LLM per key domain and keeps a target recall.}
\label{tab:intro-landscape}
\begin{tabularx}{\columnwidth}{@{}>{\raggedright\arraybackslash}p{1.4cm}YYYY>{\columncolor{sysbg}}Y@{}}
\toprule
 & \multicolumn{2}{c}{\textbf{Non-LLM}}
 & \textbf{LLM}
 & \textbf{Hybrid}
 & \textbf{Hybrid} \\
\cmidrule(lr){2-3}
 & \textbf{Auto-BI}~\cite{AutoBI}
 & \textbf{JOSIE}~\cite{DBLP:conf/sigmod/ZhuDNM19}
 & \textbf{All-pairs}\newline prompting
 & \textbf{LOTUS}~\cite{LOTUS},\newline \textbf{FDJ}~\cite{FDJ}
 & \textbf\ourmethod\newline ours \\
\midrule
\textbf{Robust to accidental overlap}
 & \xmark{} learned syntactic features
 & \xmark{} value overlap
 & \cmark{}
 & \cmark{} if routed to LLM
 & \cmark{} LLM checks each ring \\
\cmidrule{1-6}
\textbf{Finds joins with partial or zero containment}
 & \xmark{}
 & \xmark{}
 & \cmark{}
 & \xmark{} if proxy drops them
 & \cmark{} normalized formats, embeddings \\
\cmidrule{1-6}
\textbf{Labeled data needed}
 & \xmark{} corpus of BI models
 & \cmark{} none
 & \cmark{} none
 & \cmark{} none
 & \cmark{} own LLM decisions \\
\cmidrule{1-6}
\textbf{LLM calls grow with}
 & n/a
 & n/a
 & \xmark{} table pairs
 & \xmark{} candidate pairs
 & \cmark{} key domains \\
\cmidrule{1-6}
\textbf{Recall guarantee}
 & \xmark{}
 & \xmark{} top-$k$ by overlap
 & n/a (asks all pairs)
 & \cmark{} per operator
 & \cmark{} vs.\ all-pairs LLM \\
\cmidrule{1-6}
\textbf{Returns key and direction}
 & \cmark{}
 & \xmark{}
 & \cmark{}
 & \xmark{} user predicate
 & \cmark{} owner and members \\
\cmidrule{1-6}
\textbf{Consistent, non-redundant graph}
 & \cmark{} arborescence
 & \xmark{}
 & \xmark{} independent answers
 & \xmark{}
 & \cmark{} star per key domain \\
\cmidrule{1-6}
\textbf{Cost to link a new table}
 & \xmark{} re-solve global graph
 & \cmark{} index lookup
 & \xmark{} $n$ pair prompts
 & \xmark{} $n$ pair operators
 & \cmark{} affected rings only \\
\cmidrule{1-6}
% TODO(Saeed): fill from the experiments on the largest production database.
\textbf{F1; cost} (\todo{N} tables)
 & \todo{F1}; no LLM
 & \todo{F1}; no LLM
 & \todo{F1}; \todo{\$, h}
 & \todo{F1}; \todo{\$, h}
 & \todo{F1}; \todo{\$, h} \\
\bottomrule
\end{tabularx}
\end{table}


\introparagraph{Problem and Challenges}
Given a database without declared keys and access to an LLM, our goal is to discover its join keys with as few LLM calls as possible while retaining a target recall of the all-pairs LLM baseline. Table~\ref{tab:intro-landscape} situates representative systems with respect to this objective. Because tables arrive continuously, the system must also link each new table at a cost that does not grow with the size of the business. Meeting these requirements is difficult for four reasons. First, the candidate space is quadratic and almost entirely empty, so pruning must eliminate nearly all pairs before any LLM call, yet every true join that is pruned is lost silently. Second, the inexpensive signals on which pruning relies are precisely those that fail for non-LLM methods. Third, pairwise decisions ignore the structure of key domains and are therefore both redundant and potentially inconsistent. Fourth, an online system observes tables in arbitrary order: a late-arriving table may be the true owner of a key already referenced by many earlier tables, requiring earlier decisions to be revised while ensuring that the final result is independent of arrival order.

\introparagraph{Our Insight}
Enterprise joins are not distributed uniformly over table pairs. They concentrate on a modest number of \emph{key domains}--customer, project, employee, invoice identifiers, and the like. Each domain comprises one owner column that is unique and many member columns that reference it. Knowledge of the key domains immediately yields the joins: every member joins with its owner. The number of key domains is typically orders of magnitude smaller than the number of table pairs. Consequently the LLM should verify key domains rather than pairs. Independent evidence supports this premise: across 5,957 real schemas, join graphs are sparse and low-rank~\cite{Nexus}. We show that rings explain this structure, since the join graph induced by $|\mathcal{R}|$ rings has rank at most $2|\mathcal{R}|$ (Proposition~\ref{prop:rank}). Where low-rank completion leaves these factors latent, \ourmethod turns each one into an explicit object that names an entity, its owner, and the direction of every join, and that the LLM verifies as a unit. This observation lifts the linear-work insight of feature extraction (e.g., FDJ) from records to schemas and extends the value domains of LSH Ensemble with ownership and semantics: a key domain is a set of columns that refer to the same business entity, together with a designated owner. It also partitions the work according to the complementary strengths of the two classes of methods. Non-LLM signals are reliable for uniqueness and for proposing high-recall candidates, but unreliable for deciding meaning; the LLM is reliable for meaning, yet too expensive to apply to every pair.

\introparagraph{System Overview}
This paper presents \ourmethod\footnote{\ourmethod holds all columns of one key domain together, analogous to keys on a ring. Code is available to reviewers at \url{https://anonymous.4open.science/r/KeyRing}.}, a hybrid system whose LLM cost scales with the number of key domains rather than the number of table pairs. \ourmethod comprises four core components. The \emph{Column Profiler} reads the data once and materializes sketches of each column’s values, uniqueness, type, value pattern, normalized formats, and an embedding of its name and description; no subsequent step re-reads raw data. The \emph{Ring Builder} generates candidates via containment and embedding indexes tuned for recall, prunes generic and type-incompatible columns, and groups the remaining columns into \emph{rings}--one per candidate key domain--each with an elected owner. The \emph{Budgeted Verifier} decides which rings and candidates require the LLM: a local model distilled from \ourmethod’s own earlier LLM decisions accepts or rejects confident candidates; the LLM verifies each remaining ring once, in a batched prompt that presents the owner, its members, and any uncertain candidates, and removes accidental members such as surrogate overlaps. A cost-based optimizer selects thresholds and batch sizes so that \ourmethod meets a target recall $\gamma$ of the all-pairs LLM baseline with probability at least $1-\delta$, accounting for the correlated errors that arise when a single verdict decides an entire ring. Confirmed rings propagate joins by transitivity without further LLM calls. The output is a star-shaped graph centered on each owner rather than a clique of redundant edges.

\ourmethod operates in two modes over the same persistent state (profiles, indexes, rings, and the distilled model). In \emph{bootstrap} mode it constructs the rings of an entire business from all of its tables. In \emph{online} mode it links a single arriving table: it profiles only the new columns, probes the existing indexes, and assigns each key-like column to an existing ring or initiates a new one. The \emph{Ring Maintainer} revises only the affected rings--re-electing a superior owner, merging rings that the new table connects, or splitting a ring that verification shows to mix distinct domains--and the Budgeted Verifier re-checks solely those rings. Each change updates the edges of a single ring; consequently the cost of an arrival grows with the size of the new table and of the rings it touches, not with the total number of tables. Bootstrap is simply the special case in which all tables arrive simultaneously; we prove that the final rings are independent of arrival order.

\introparagraph{Contributions}
Our primary contribution is \ourmethod, a practical hybrid system for join-key discovery over thousands of tables that is more accurate than non-LLM methods and whose LLM cost scales with the number of key domains rather than the number of table pairs. The concrete technical contributions are:
\begin{itemize}[leftmargin=*]
\item An analysis of the failure modes of non-LLM, LLM-based, and hybrid join discovery on enterprise data; a formalization of cost-bounded join-key discovery with a recall target relative to an LLM oracle; a structural result showing that the join graph induced by rings has rank at most twice the number of rings, which explains the low-rank structure of real join graphs and turns its latent factors into explicit, verifiable key domains with owners and directions; and a system architecture that solves the problem by discovering key domains instead of judging table pairs (Sections~\ref{sec:background} and~\ref{sec:architecture}).
\item A ring builder that converts a single pass of column profiling into candidate rings with elected owners, employing containment and embedding indexes together with format normalization, and that prunes generic columns that match too many owners--all without LLM calls and without materializing the quadratic space of column pairs, so that it scales to tens of thousands of columns per business (Section~\ref{sec:ring_builder}).
\item A budgeted verifier that routes each candidate to a distilled local model or to one batched LLM prompt per ring, together with a cost-based optimizer whose thresholds guarantee a target recall of the all-pairs LLM baseline with prescribed confidence, while correctly accounting for correlated ring-level errors (Section~\ref{sec:verifier}).
\item A ring maintainer for the online setting that identifies the minimal set of rings affected by an arriving table, re-verifies only those rings, and performs owner re-election, merges, and splits, so that arrival cost depends on the new table and the rings it touches rather than on the total number of tables, and the final result is independent of arrival order (Section~\ref{sec:online}).
\item A comprehensive experimental evaluation on \todo{N} production enterprise databases containing up to \todo{8K} tables and on public benchmarks, comparing against non-LLM methods (inclusion-dependency discovery, Auto-BI, joinable-table search), all-pairs LLM prompting, hybrid baselines, and global-structure inference with low-rank completion (Nexus) (Section~\ref{sec:experiments}). Under the observed operating regime \ourmethod improves join F1 over the strongest non-LLM baseline, reduces LLM calls, tokens, and end-to-end time by a factor of 2--5$\times$ relative to all-pairs prompting while retaining the large majority of its joins, and attains high precision on human-labeled joins. When tables are replayed in creation order, each arriving table is linked with substantially fewer LLM calls than pairwise verification and the final graph coincides with that of a full bootstrap run. Downstream analytic and Text-to-SQL workloads that consume the resulting compact graph exhibit fewer intermediate results and lower execution cost without altering query answers.
\end{itemize}
